\subsection{Identidad de Suma de Fibonacci + Fast Doubling}

\graybold{Preguntas guía:}

\begin{itemize}
\item ¿Necesito un término de Fibonacci (o una suma de varios) con índices enormes (hasta $10^9$ o más), descartando cualquier enfoque $O(n)$?
\item ¿La pregunta involucra una SUMA de un rango de términos de Fibonacci, no solo un término aislado?
\item ¿Hay identidades cerradas conocidas que reduzcan la suma a un par de evaluaciones puntuales de Fibonacci?
\end{itemize}

\graybold{¿Para qué sirve?} Calcular términos de Fibonacci con índices astronómicos en tiempo logarítmico (fast doubling), y combinarlo con identidades de suma para responder consultas de rango sin iterar término por término.

\graybold{Complejidad:} \bigO{log(max(N,M))} por consulta.

\graybold{Fundamento teórico:} La suma de prefijo $S(n) = \sum_{i=0}^{n} F(i) = F(n+2) - 1$ es una identidad clásica (demostrable por inducción). De ahí, $\sum_{i=N}^{M} F(i) = S(M) - S(N-1) = F(M+2) - F(N+1)$, válida incluso en el borde $N=0$ (ya que $S(-1)=F(1)-1=0$). Para evaluar $F(k) \bmod p$ con $k$ grande, se usa \emph{fast doubling}: las identidades $F(2n) = F(n)\cdot(2F(n+1)-F(n))$ y $F(2n+1) = F(n)^2 + F(n+1)^2$ permiten calcular el par $(F(n), F(n+1))$ recursivamente en $O(\log n)$, dividiendo el índice a la mitad en cada llamada (análogo a exponenciación rápida).

\graybold{Ejercicio aplicado}

Dados N y M, calcular $(F(N) + F(N+1) + \cdots + F(M)) \bmod 10^9+7$.

\graybold{I/O:} T ≤ 1000, N,M ≤ 10\textsuperscript{9} → lectura total con tokenización (patrón 2); cada consulta se resuelve con 2 llamadas a fast doubling.

\graybold{Código solución}

\begin{lstlisting}[language=Python]
import sys

MOD = 1000000007

def fib_pair(n):
    if n == 0:
        return (0, 1)
    a, b = fib_pair(n >> 1)
    c = (a * ((2 * b - a) % MOD)) % MOD
    d = (a * a + b * b) % MOD
    if n & 1:
        return (d, (c + d) % MOD)
    return (c, d)

def fib(n):
    return fib_pair(n)[0]

def solve():
    data = sys.stdin.buffer.read().split()
    t = int(data[0])
    out = []
    idx = 1
    for _ in range(t):
        n = int(data[idx]); m = int(data[idx + 1]); idx += 2
        # suma(N..M) = F(M+2) - F(N+1)  (usando S(n) = suma(0..n) = F(n+2) - 1)
        ans = (fib(m + 2) - fib(n + 1)) % MOD
        out.append(str(ans))
    sys.stdout.write("\n".join(out) + "\n")

solve()
\end{lstlisting}
